% Source: https://artofproblemsolving.com/community/c3h1076780_max_and_min
% Stable identity: aops-topic-1076780-q01
The denominator $2x^2+2x+1=2\left(x+\dfrac{1}{2}\right)^2+\dfrac{1}{2}$ is always positive, so $y$ is defined for all real $x$. A value $y$ is attained exactly when
\[
(2y-1)x^2+(2y+2)x+(y+3)=0
\]
has a real solution $x$. If $y=\dfrac{1}{2}$ this is $3x+\dfrac{7}{2}=0$, which has a solution. Otherwise we need the discriminant to be nonnegative:
\[
(2y+2)^2-4(2y-1)(y+3)=-4y^2-12y+16=-4(y+4)(y-1)\ge0\iff-4\le y\le1 .
\]
Both ends are attained: $y=1$ at $x=-2$ (value $\dfrac{5}{5}$) and $y=-4$ at $x=-\dfrac{1}{3}$ (value $\dfrac{-\dfrac{20}{9}}{\dfrac{5}{9}}$). So the maximum is $1$, the minimum is $-4$, and their sum is $\boxed{-3}$.
